\section{The XRay--BOLT Interaction on AArch64}
\label{sec:xray}

Running BOLT on an AArch64 ClickHouse-derived binary has a reputation for
fragility. We present a root-cause analysis that reframes the problem: the
failure we observed is not, at bottom, a defect of \texttt{llvm-bolt}, but a
consequence of an instrumentation feature that a release binary should not
carry in the first place.

\subsection{Symptom}

On AArch64, BOLT-processed binaries failed at run time. In our deployment the
failure surfaced indirectly: ephemeral CI runners launched from such binaries
reported \texttt{VssOAuthTokenRequestException: The signature is not valid}
during the broker handshake and never came online --- a misleading symptom
several layers removed from the true cause. Direct BOLT invocation, when forced,
failed earlier with \texttt{error: expected relocatable expression} during
emission, and an assertion build of \texttt{llvm-bolt} (v22.1.1) pinpointed the
failure to
\begin{center}
\texttt{Relocation::getComposeOpcodeFor():}\\
\texttt{assert(Arch == Triple::riscv64 \&\& "only implemented for RISC-V")}
\end{center}

\subsection{Root cause}

The function that BOLT uses to \emph{compose} paired relocations into a single
MC expression is implemented only for RISC-V. On AArch64 it is reached by
composed \texttt{R\_AARCH64\_PREL64} relocations --- PC-relative 64-bit data
references of the form $S + A - P$. Instrumenting the binary revealed that the
overwhelming majority of these composed groups, including one pathological group
of \textbf{$62{,}692$ relocations sharing a single offset}, lived in the
\texttt{xray\_instr\_map} section. In other words, the relocations BOLT could
not lower were emitted by \textbf{XRay instrumentation}, which was enabled by
default.

XRay~\cite{xray} is a function-call tracing facility: it inserts entry/exit
``sleds'' at every function and a side table (\texttt{xray\_instr\_map},
\texttt{xray\_fn\_idx}) of PC-relative pointers into them. We measured these
sections at $\sim\!45$\,MB in the datastore binary. A release database server
has no business carrying always-on tracing sleds; the feature is appropriate
for a profiling build, opt-in, and nowhere else.

\subsection{Fix}

The defect upstream was a default: \texttt{ENABLE\_XRAY} was \texttt{ON}.
Turning it off (with the coverage build that legitimately needs XRay
re-enabling it locally) has three effects:

\begin{enumerate}[leftmargin=1.5em]
  \item The \texttt{xray\_instr\_map} section disappears, taking the
  unlowerable composed \texttt{R\_AARCH64\_PREL64} relocations with it. BOLT
  then emits cleanly on AArch64 and the resulting binary runs correctly.
  \item The binary shrinks from $945$\,MB to $648$\,MB after BOLT --- a $31\%$
  reduction --- before any other consideration.
  \item Per-function sled overhead is removed from the hot path.
\end{enumerate}

A residual, genuinely-upstreamable robustness fix remains --- teaching
\texttt{getComposeOpcodeFor} the AArch64 cases so that \emph{any} future
composed \texttt{PREL64} (not only XRay's) lowers correctly --- but it is not
required for Datastore, because Datastore should not emit those relocations at
all. Two operational notes for reproducing BOLT on the large AArch64 binary:
\texttt{llvm-bolt} recurses deeply during veneer insertion and overflows the
default $8$\,MB stack, so it must run under a raised stack limit
(\texttt{ulimit -s 2{,}097{,}152}); and the GB10 firmware we tested does not
expose the SPE table, so branch profiles are collected with cycle sampling and
\texttt{perf2bolt --nl} rather than via LBR/SPE.
